\documentclass[12pt]{article}

\usepackage[spanish]{babel}
\usepackage[utf8]{inputenc}
\usepackage[T1]{fontenc}
\usepackage{graphicx}
\usepackage{geometry}
\usepackage{listings}
\usepackage{xcolor}
\usepackage{float}
\usepackage{array}
\usepackage{hyperref}

\geometry{a4paper, margin=2.5cm}

\title{
Proyecto Perritos \\ 
Aplicación de Registro e Interacción de Mascotas
}

\author{
Miranda Rodriguez Cruz
}

\date{2026}

\definecolor{codegreen}{rgb}{0,0.6,0}
\definecolor{codegray}{rgb}{0.5,0.5,0.5}
\definecolor{codepurple}{rgb}{0.58,0,0.82}
\definecolor{backcolour}{rgb}{0.95,0.95,0.92}

\lstdefinestyle{mystyle}{
    backgroundcolor=\color{backcolour},
    commentstyle=\color{codegreen},
    keywordstyle=\color{blue},
    numberstyle=\tiny\color{codegray},
    stringstyle=\color{codepurple},
    basicstyle=\ttfamily\small,
    breaklines=true,
    captionpos=b,
    keepspaces=true,
    numbers=left,
    numbersep=5pt,
    showspaces=false,
    showstringspaces=false,
    showtabs=false,
    tabsize=2
}

\lstset{style=mystyle}

\begin{document}

\maketitle

\tableofcontents

\newpage

% ---------------------------------------------------
% INTRODUCCION
% ---------------------------------------------------

\section{Introducción}

En este proyecto desarrollé una aplicación llamada \textbf{Perritos}, utilizando Flutter y Dart, con el objetivo de crear una pequeña red social enfocada en mascotas.

La aplicación permite registrar perritos, agregar imágenes, comentar publicaciones y reaccionar mediante diferentes emojis. Además, integré Cloudinary para almacenar imágenes en la nube y Render para desplegar la aplicación web.

Durante el desarrollo del proyecto me enfoqué principalmente en mantener una estructura organizada y reutilizable mediante el uso de patrones de diseño y principios SOLID.

Una de las partes más importantes del proyecto fue separar correctamente responsabilidades para evitar código repetido y facilitar el mantenimiento del sistema.

% ---------------------------------------------------
% OBJETIVOS
% ---------------------------------------------------

\section{Objetivo General}

Desarrollar una aplicación multiplataforma que permita registrar mascotas e interactuar mediante comentarios y reacciones utilizando tecnologías modernas y servicios en la nube.

\section{Objetivos Específicos}

\begin{itemize}
    \item Implementar autenticación de usuarios.
    \item Registrar mascotas dentro de la aplicación.
    \item Permitir comentarios y reacciones.
    \item Subir imágenes mediante Cloudinary.
    \item Aplicar patrones de diseño.
    \item Aplicar principios SOLID.
    \item Desplegar la aplicación utilizando Render.
    \item Generar un APK funcional.
\end{itemize}

% ---------------------------------------------------
% TECNOLOGIAS
% ---------------------------------------------------

\section{Tecnologías Utilizadas}

\begin{itemize}
    \item Flutter
    \item Dart
    \item Cloudinary
    \item Render
    \item GitHub
    \item Visual Studio Code
    \item Overleaf
\end{itemize}

% ---------------------------------------------------
% CAPTURAS
% ---------------------------------------------------

\section{Capturas del Sistema}

\subsection{Pantalla de Login}

\begin{figure}[H]
\centering
\includegraphics[width=0.9\textwidth]{login.png}
\caption{Pantalla de inicio de sesión}
\end{figure}

\subsection{Registro de Mascotas}

\begin{figure}[H]
\centering
\includegraphics[width=0.9\textwidth]{registro.png}
\caption{Formulario de registro de mascotas}
\end{figure}

\subsection{Comentarios y Reacciones}

\begin{figure}[H]
\centering
\includegraphics[width=0.9\textwidth]{comentarios.png}
\caption{Sistema de comentarios y reacciones}
\end{figure}

% ---------------------------------------------------
% SOLID
% ---------------------------------------------------

\section{Principios SOLID Aplicados}

Durante el desarrollo del proyecto apliqué algunos principios SOLID con el objetivo de mantener el código más organizado, reutilizable y fácil de mantener.

\begin{table}[H]
\centering
\begin{tabular}{|p{4cm}|p{9cm}|}
\hline

\textbf{Principio} & \textbf{Aplicación en el Proyecto} \\

\hline

Single Responsibility Principle (SRP)
&
Cada clase fue creada con una responsabilidad específica. \\

\hline

Open Closed Principle (OCP)
&
El sistema puede extenderse agregando nuevas funcionalidades sin modificar completamente la estructura principal. \\

\hline

Liskov Substitution Principle (LSP)
&
Los widgets reutilizables pueden sustituirse manteniendo compatibilidad visual y funcional. \\

\hline

Interface Segregation Principle (ISP)
&
Los servicios fueron separados según sus responsabilidades. \\

\hline

Dependency Inversion Principle (DIP)
&
Las pantallas dependen de servicios especializados para manejar funcionalidades externas. \\

\hline

\end{tabular}

\caption{Principios SOLID utilizados}
\end{table}

% ---------------------------------------------------
% SRP
% ---------------------------------------------------

\subsection{Single Responsibility Principle (SRP)}

Este principio establece que una clase debe tener únicamente una responsabilidad.

En mi proyecto decidí separar completamente la lógica relacionada con Cloudinary dentro de un servicio independiente.

\begin{figure}[H]
\centering
\includegraphics[width=1\textwidth]{cloudinary.png}
\caption{Servicio Cloudinary}
\end{figure}

\begin{lstlisting}[language=Dart]
class CloudinaryService {

  static Future<String?> uploadImageWeb(
    Uint8List imageBytes,
  ) async {

  }
}
\end{lstlisting}

\textbf{Explicación}

La clase CloudinaryService únicamente se encarga de subir imágenes a Cloudinary.

No contiene lógica relacionada con comentarios, mascotas o interfaz gráfica.

Esto me permitió mantener una mejor organización y evitar mezclar responsabilidades.

% ---------------------------------------------------
% OCP
% ---------------------------------------------------

\subsection{Open Closed Principle (OCP)}

Este principio establece que el sistema debe estar abierto para extensión pero cerrado para modificación.

En el proyecto implementé las reacciones utilizando botones independientes, lo cual permite agregar nuevas reacciones fácilmente sin modificar completamente la lógica principal.

\begin{lstlisting}[language=Dart]
Wrap(
  spacing: 10,
  children: [

    ElevatedButton(
      child: Text('❤️'),
    ),

    ElevatedButton(
      child: Text('��'),
    ),

    ElevatedButton(
      child: Text('��'),
    ),
  ],
)
\end{lstlisting}

\textbf{Explicación}

Gracias a esta estructura puedo agregar nuevas reacciones posteriormente sin afectar otras partes del sistema.

Esto facilita la escalabilidad y mantenimiento de la aplicación.

% ---------------------------------------------------
% LSP
% ---------------------------------------------------

\subsection{Liskov Substitution Principle (LSP)}

Este principio establece que los componentes deben poder sustituirse sin afectar el funcionamiento general del sistema.

Flutter trabaja naturalmente mediante widgets reutilizables.

\begin{lstlisting}[language=Dart]
Widget _tarjetaCampo(
  String titulo,
  TextEditingController controller,
)
\end{lstlisting}

\textbf{Explicación}

En este caso el widget puede reutilizarse múltiples veces dentro de la interfaz manteniendo el mismo comportamiento y estructura visual.

Esto me permitió reducir código repetido y mantener consistencia en el diseño.

% ---------------------------------------------------
% ISP
% ---------------------------------------------------

\subsection{Interface Segregation Principle (ISP)}

Este principio establece que las clases no deben depender de funcionalidades que no utilizan.

En el proyecto decidí separar los servicios según sus funciones específicas.

\begin{lstlisting}[language=Dart]
final _auth = AuthService();

final _picker = ImagePicker();
\end{lstlisting}

\textbf{Explicación}

Cada servicio tiene una responsabilidad independiente.

AuthService se encarga de autenticación mientras que ImagePicker únicamente maneja selección de imágenes.

Esto mejora la modularidad y organización del proyecto.

% ---------------------------------------------------
% DIP
% ---------------------------------------------------

\subsection{Dependency Inversion Principle (DIP)}

Este principio establece que las clases principales deben depender de abstracciones y servicios especializados.

\begin{lstlisting}[language=Dart]
final imageUrl =
  await CloudinaryService.uploadImageWeb(
    imageBytes,
  );
\end{lstlisting}

\textbf{Explicación}

La pantalla principal no implementa directamente toda la lógica HTTP para subir imágenes.

En su lugar, delega esa responsabilidad a CloudinaryService.

Esto facilita reutilizar el servicio y mantener una mejor separación de responsabilidades.

% ---------------------------------------------------
% PATRONES
% ---------------------------------------------------

\section{Patrones de Diseño Utilizados}

Durante el desarrollo del proyecto utilicé diferentes patrones de diseño con el objetivo de mantener una estructura organizada, modular y fácil de mantener.

Cada patrón fue utilizado para resolver problemas específicos dentro de la aplicación.

% ---------------------------------------------------
% MVC
% ---------------------------------------------------

\subsection{Patrón MVC}

El patrón MVC fue uno de los patrones principales utilizados en el proyecto.

MVC significa:

\begin{itemize}
    \item Model
    \item View
    \item Controller
\end{itemize}

La finalidad de este patrón es separar los datos, la interfaz gráfica y la lógica de control.

Esto ayuda a mantener una estructura mucho más organizada.

% ---------------------------------------------------
% MODEL
% ---------------------------------------------------

\subsubsection{Model}

El Model representa los datos de la aplicación.

\begin{figure}[H]
\centering
\includegraphics[width=1\textwidth]{modelo.png}
\caption{Modelo Dog y Comentario}
\end{figure}

\begin{lstlisting}[language=Dart]
class Dog {

  final String nombre;
  final String raza;
  final String edad;
  final String color;

  List<Comentario> comentarios;
}
\end{lstlisting}

\textbf{Explicación}

La clase Dog representa la estructura principal de datos utilizada para almacenar información de cada mascota.

Además, contiene una lista de comentarios relacionados con cada publicación.

% ---------------------------------------------------
% VIEW
% ---------------------------------------------------

\subsubsection{View}

La View representa la interfaz gráfica mostrada al usuario.

\begin{lstlisting}[language=Dart]
class PerritosScreen
extends StatefulWidget {

  const PerritosScreen({
    super.key
  });
}
\end{lstlisting}

\textbf{Explicación}

PerritosScreen representa la pantalla principal donde el usuario puede visualizar publicaciones, comentarios y reacciones.

% ---------------------------------------------------
% CONTROLLER
% ---------------------------------------------------

\subsubsection{Controller}

El Controller se encarga de manejar eventos y actualizar la aplicación.

\begin{figure}[H]
\centering
\includegraphics[width=1\textwidth]{perritos_screen.png}
\caption{Lógica de PerritosScreen}
\end{figure}

\begin{lstlisting}[language=Dart]
Future<void> _guardar() async {

  final nuevoPerro = Dog(
    nombre: nombre,
    raza: raza,
    edad: edad,
    color: color,
    imageUrl: imageUrl,
  );

  setState(() {
    perros.add(nuevoPerro);
  });
}
\end{lstlisting}

\textbf{Explicación}

La función _guardar() controla el proceso de registro de mascotas.

Aquí se crean objetos Dog y posteriormente se actualiza el estado de la interfaz utilizando setState().

% ---------------------------------------------------
% COMPONENT BASED
% ---------------------------------------------------

\subsection{Component-Based Architecture}

Flutter trabaja utilizando componentes reutilizables llamados widgets.

Durante el proyecto utilicé componentes personalizados para evitar repetir código.

\begin{lstlisting}[language=Dart]
Widget _tarjetaCampo(
  String titulo,
  TextEditingController controller,
)
\end{lstlisting}

\textbf{Explicación}

Este componente fue utilizado múltiples veces dentro del formulario principal.

Gracias a esto pude mantener una interfaz más consistente y reducir código repetido.

% ---------------------------------------------------
% STATE MANAGEMENT
% ---------------------------------------------------

\subsection{State Management}

Para manejar cambios dinámicos dentro de la aplicación utilicé setState().

\begin{lstlisting}[language=Dart]
setState(() {
  perros.add(nuevoPerro);
});
\end{lstlisting}

\textbf{Explicación}

setState() actualiza automáticamente la interfaz cuando se agregan nuevas mascotas, comentarios o reacciones.

Esto permitió mantener una interacción dinámica dentro de la aplicación.

% ---------------------------------------------------
% SERVICE PATTERN
% ---------------------------------------------------

\subsection{Service Pattern}

También utilicé el patrón Service para separar funcionalidades externas dentro de servicios independientes.

\begin{lstlisting}[language=Dart]
class CloudinaryService {

}
\end{lstlisting}

\textbf{Explicación}

CloudinaryService encapsula toda la lógica relacionada con Cloudinary.

Esto permitió reutilizar fácilmente el servicio y mantener una mejor organización.

% ---------------------------------------------------
% SINGLETON
% ---------------------------------------------------

\subsection{Singleton}

El patrón Singleton busca garantizar que exista una única instancia compartida de un servicio.

\begin{lstlisting}[language=Dart]
final _auth = AuthService();
\end{lstlisting}

\textbf{Explicación}

En el proyecto reutilicé una única instancia de AuthService para manejar autenticación dentro de diferentes pantallas.

Esto evita crear múltiples instancias innecesarias y mejora la administración global del servicio.

% ---------------------------------------------------
% CLOUDINARY
% ---------------------------------------------------

\section{Integración con Cloudinary}

Cloudinary fue utilizado para almacenar imágenes en la nube.

El flujo implementado fue el siguiente:

\begin{enumerate}
    \item El usuario selecciona una imagen.
    \item Flutter envía la imagen a Cloudinary.
    \item Cloudinary devuelve una URL.
    \item La URL se guarda dentro del objeto Dog.
\end{enumerate}

\begin{lstlisting}[language=Dart]
final imageUrl =
  await CloudinaryService.uploadImageWeb(
    imageBytes,
  );
\end{lstlisting}

% ---------------------------------------------------
% RENDER
% ---------------------------------------------------

\section{Despliegue en Render}

La aplicación fue desplegada utilizando Render como plataforma de hosting web.

Para realizar el despliegue primero actualicé el proyecto en GitHub y posteriormente configuré un Static Site dentro de Render.

% ---------------------------------------------------
% APK
% ---------------------------------------------------

\section{Generación APK}

Finalmente generé un APK funcional utilizando Flutter.

\begin{lstlisting}[language=bash]
flutter build apk
\end{lstlisting}

% ---------------------------------------------------
% CONCLUSIONES
% ---------------------------------------------------

\section{Conclusiones}

Durante el desarrollo de este proyecto comprendí la importancia de mantener una estructura organizada y modular.

La implementación de patrones de diseño y principios SOLID me permitió separar responsabilidades, reutilizar código y facilitar el mantenimiento del sistema.

Además, integrar tecnologías como Cloudinary y Render permitió acercar el proyecto a un entorno más realista y profesional.

% ---------------------------------------------------
% REFERENCIAS
% ---------------------------------------------------

\section{Referencias}

\begin{itemize}
    \item Flutter Documentation
    \item Dart Documentation
    \item Cloudinary Documentation
    \item Render Documentation
    \item SOLID Principles
\end{itemize}

\end{document}